\subsection{A New Event During Recovery in a Four-Family Model}
\label{app:four-family-conditional}

We tested whether the response still visible after a fault clears or a
line recloses can account for an apparent second event. This experiment
uses a reduced IEEE 39-bus simulator with classical machine dynamics,
governors and eight PMUs. Its four interventions are a bus fault, a
temporary line opening, a positive pulse in generator torque and a load
change. Each record lasts 24~s, including a one-second quiet reference.
Both the simulator and the predictor use the same network and machine
equations. The experiment tests uncertainty in their parameters.
Li and Wang predict and subtract a preceding disturbance to classify a
second event from a triggered PMU window \cite{liwang2019successive}.
Here the physical fit supplies the competing prediction that recovery alone
continues. The monitor must decide whether to issue a second alarm and,
if it does, name the affected bus or line.

\paragraph{A common prediction after clearance.}
The first declaration supplies the family and asset of a fault or line
opening. When the inferred probability that clearance occurred at least
15 reports earlier reaches 0.99, we refine the switch times from nearby
voltage changes. Let $k_f$ be the fitting report, $y_j$ the complex PMU
voltages normalized to their quiet mean, and $V_j(\theta)$ the predicted
voltages for the inferred switch history. Both monitors fit
\[
 \widehat\theta=\arg\min_{\theta\in\Theta}
 \sum_{j=\widehat\tau}^{k_f}
 \frac{\|y_j-V_j(\theta)\|_2^2}{(0.001)^2}
 +\sum_{i=1}^{4}\frac{\log^2(\theta_i/\theta_{0,i})}{s_i^2}.
\]
The four penalized parameters are inertia and damping multipliers,
governor time constant and droop, with nominal values $(1,1,0.5,0.05)$
and log scales $(0.25,0.65,0.5,0.45)$. For a fault, the fit includes a
fifth bounded parameter for conductance. With the asset back in service,
the fitted swing and governor dynamics predict its residual motion.
The network equations map that motion to the eight PMU voltages and
oriented branch currents. The reference monitor retains its original bank of new-event hypotheses
and can also declare a generator event from the governor-fitted recovery
residual. The original stream has priority at
each publication time.

The alternative uses one decision after the fit becomes available.
It subtracts the predicted departure due to the first event from each
new report and compares all four intervention families with continued
recovery in those coordinates. Both hypotheses then account for the
same first event. A conditional filter initialized from the quiet
reference starts with the next report. The measurements used to fit the
first event do not enter its subsequent event likelihood again.

For a candidate with predicted innovation $s$, let $r$ be the observed
innovation in the 31 whitened measurement coordinates, and $v$ the
baseline predictive scale. Under a multivariate Student law with
$\nu$ degrees of freedom and common scale matrix $c_b v I$, the
single-report log likelihood ratio against continued recovery, before a
second event has been declared, is
\[
 \frac{\nu+31}{2}
 \log\!\left(\frac{\nu c_b v+\|r\|^2}
                   {\nu c_b v+\|r-s\|^2}\right).
\]
The filter adds these log ratios across reports for each candidate history,
then sums the prior-weighted likelihoods over onset, asset and duration.
A previously fitted slope and intercept calibrate the resulting log odds.
Every three reports the monitor compares the expected cost of publishing
each family--asset pair with the cost of waiting three more reports. It
approximates that one-step expectation by sampling predictive measurements.
Under the null model, $\|r\|^2/(31c_b v)$ follows $F_{31,\nu}$.
Here $\nu=3$. We estimated $c_b$ by dividing the median observed radial energy
$\|r\|^2/(31v)$ by the median of this reference distribution.
Twenty isolated training faults and sixteen isolated line openings supplied the residuals.
Their labeled switches initialized the training fits. In operation the
monitor infers the switches. We fitted one scale at each of three
recovery ages (up to 45, 46--150, and beyond 150 reports after
clearance). The fitted scale multipliers were 0.866, 0.856 and 0.840,
respectively. The unit-scale control uses the same conditional filter.
Conditioning on the fitted trajectory is a plug-in approximation to
integration over uncertain physical parameters.

We checked whether these parameters leave an asset-specific pattern in
the prediction error. In 10 isolated faults and 11 isolated line openings
from development, we stacked the real and imaginary voltage errors from
eight PMUs over 15 consecutive reports. We projected each 240-coordinate
vector onto the span of the voltage sensitivities to the fitted parameters.
Across 728 nonoverlapping windows after fitting, the median projected
energy fraction (first within each record, then across records) was
0.973 for faults and 0.617 for lines. Cyclically shifting the PMU
channels gave 0.965 and 0.616. The high alignment survives a change of
sensor labels, which limits its value as evidence for the location.

\paragraph{What duration can identify.}
The original generator prior admitted durations of 27--66 reports.
We assigned half of each generator's prior mass to pulses of 6--12
reports, preserving the total prior by asset and the original torque
quadrature. At a fixed onset, asset and torque, all durations still
active at report $k$ have the same input history and the same prefix
likelihood $L_{\rm active}$. Their contribution to the evidence is
$L_{\rm active}\Pr(D>k-\tau)$ for onset report $\tau$ and duration $D$.
A duration change can alter the relative weight of already-cleared
paths. Distinguishing two pulses that have both remained active requires
a later report at which their trajectories differ.

\paragraph{Paired evaluation.}
Each scenario block shares its noise, first-event realization and
timing across an isolated first event and its two-event counterparts.
First- and second-event strengths are drawn independently.
Each block tests nominal regulation and a second setting with time constant and droop drawn
independently and uniformly over 0.35--1.05~s and 0.035--0.085,
respectively. The supervised
comparator uses the source-excluded ExtraTrees refit with its original
decision rules. All methods consume the same reports in time order.
We froze the estimators and operating points before generating these
records. Event times, disturbances and noise were drawn afresh. The
reserved physical assets were the same as in our earlier four-family
experiments, so the new trajectories do not expand the set of tested
locations.
A successful joint decision requires the correct family and exact
bus or line within 45 reports of onset. A second count requires
publication before the event ends. Macro-F1 weights the four families
equally. Extras include every unmatched alarm in event and quiet records.

\begin{table*}[t]\centering\footnotesize
\caption{Paired four-family stress test on the reduced IEEE 39-bus plant. The 420 events are the denominator for family and joint accuracy. Joint requires the correct asset and family within 1.5 s. Active also requires publication before clearance. Extras are unmatched alarms on all records.}
\label{tab:four-family-conditional}\setlength{\tabcolsep}{5pt}
\begin{tabular}{lrrrrrr}\toprule
Decision model & Detected & Family & Joint & Active & Macro-F1 & Extras\\\midrule
Reference & 403 & 403 & 397 & 391 & 0.9145 & 58 \\
Short torque pulses & 411 & 411 & 405 & 391 & 0.8848 & 117 \\
Conditional recovery & 409 & 409 & 407 & 400 & 0.9444 & 38 \\
Conditional recovery and short pulses & 417 & 417 & 415 & 400 & 0.9502 & 41 \\
Combined, unit scale & 417 & 417 & 415 & 400 & 0.9502 & 41 \\
Refitted ExtraTrees & 294 & 252 & 240 & 63 & 0.6786 & 2 \\
\midrule\multicolumn{7}{l}{Operating points selected on separate calibration blocks}\\
Reference ($b=16$) & 395 & 395 & 391 & 390 & 0.9430 & 21 \\
Conditional recovery and short pulses ($b=16$) & 410 & 410 & 408 & 396 & 0.9809 & 7 \\
\bottomrule\end{tabular}\end{table*}

The blind trial contained 465 records in 15 independent scenario blocks,
with 165 distinct quiet realizations (1.05 monitored hours).
The second-event generator buses (32, 35 and 38) and all other test assets were excluded from labeled model fitting and calibration.
Table~\ref{tab:four-family-conditional} reports a factorial comparison at the reference operating point.
As a secondary comparison, six separate calibration blocks set the alarm budget to the number of extras from the frozen trees.
For the reference and combined monitors we divided fitted event odds by $b\in\{1,4,16\}$ and selected the most accurate joint decision satisfying the budget.
When all three violated the budget, we selected the smallest excess before comparing accuracy.
The selected reference rule gave 0 calibration extras and 151 joint identifications out of 168.
The selected combined rule gave 2 calibration extras and 154 joint identifications out of 168.
The trees produced 0 calibration extras.

At the original odds, the combined rule changed joint accuracy by +4.29
percentage points (paired 95\% block-bootstrap interval [1.43, 7.14]),
macro-F1 by +3.57 points (interval [-0.24, 7.52]),
and unmatched alarms by -17.
With the same physical filter and unit scale the joint count was 415 and there were 41 extras.
The combined rule produced 0 alarms in quiet records and 41 unmatched alarms in event records.
The conditional forecast raised joint diagnoses from 25/30 to 30/30 both for loads after faults and for second line openings.
Extending the duration prior raised joint diagnoses for isolated short generator pulses from 20/30 to 28/30. In both cases only 14/30 were identified before pulse removal.
The combined rule and the reference identified the same number of generator events following a fault (28/30) or a line opening (29/30).
Across all 30 generator--amplitude combinations, the six-report
innovation prefix of the shortest pulse was identical to that of the corresponding long pulse in the implemented bank.

At the operating points selected on calibration, the combined model identified 408 family--asset pairs against 391 for the reference.
The paired gains were 4.05 percentage points in joint accuracy (95\% interval [2.14, 6.19]) and 3.79 points in macro-F1 ([1.60, 6.36]).
The combined model exceeded 0.96 in macro-F1, family accuracy and joint accuracy within the deadline. Of the 420 events, 396 had correct joint decisions before clearance.
It issued seven extras against two from the trees. Three followed isolated faults. In the records with two line openings, three were duplicate diagnoses during the second opening and one followed reclosure.
Thus the combined model missed the prespecified macro-F1 target at the reference odds and exceeded the tree alarm budget at both operating points.
These comparisons use simulated positive torque pulses, a fixed PMU installation and known topology.
The resulting operating points need independent validation on a different plant and a longer continuous stream.

